% Exemplo: primeira aula do curso (subfiles). O preâmbulo vem de ../curso.tex.
% Compilar só esta aula: latexmk -pdf aula-01.tex   Testado com: beamer, subfiles, TeX Live 2024+
%%BEGIN:aulainicio
\documentclass[../curso.tex]{subfiles}

\begin{document}

\part{O que é o \LaTeX}
\begin{refsection}       % referências próprias desta aula
%%END:aulainicio

\section{Motivação}

%%BEGIN:aulablocos
\begin{frame}{Texto e formatação separados}
  \begin{block}{Definição}
    O \LaTeX\ é um sistema de preparação de documentos: o autor escreve o conteúdo e a
    estrutura, e o programa cuida da aparência \parencite{lamport1994}.
  \end{block}
  \begin{exampleblock}{Exemplo}
    \texttt{\textbackslash section\{Introdução\}} cria uma seção numerada e a inclui no sumário.
  \end{exampleblock}
  \begin{alertblock}{Atenção}
    O arquivo \texttt{.tex} precisa ser compilado para gerar o PDF.
  \end{alertblock}
\end{frame}
%%END:aulablocos

\begin{frame}{Vantagens}
  \begin{itemize}[<+->]
    \item Qualidade tipográfica, sobretudo em fórmulas
    \item Numeração e referências automáticas
    \item Arquivos de texto, fáceis de versionar
  \end{itemize}
\end{frame}

\section{Referências}

\begin{frame}[allowframebreaks]{Referências}
  \printbibliography[heading=none]
\end{frame}

\end{refsection}
\end{document}
